% MIMN：新行为经控制器寻址固定记忆槽，读出并擦除写入；候选只在CTR读出阶段进入。
% 单栏图（figure），内容宽约80mm；上行为新行为→控制器，记忆矩阵在下方，读出标注独立置于框下。
\begin{tikzpicture}[x=1mm,y=1mm,
  font=\figurefont\footnotesize,
  mbeh/.style={rectangle,draw=FigureInput!70,fill=FigureInput!12,align=center,
    text=BookInk,rounded corners=1.4pt,minimum height=9mm,text width=18mm,inner sep=2pt},
  mctl/.style={rectangle,draw=FigureModel!70,thick,fill=FigureModel!10,align=center,
    text=BookInk,rounded corners=1.4pt,minimum height=16mm,text width=24mm,inner sep=2pt},
  mmem/.style={rectangle,draw=FigureLoss!75,fill=FigureLoss!12,align=center,
    text=BookInk,rounded corners=1.4pt,minimum height=26mm,text width=18mm,inner sep=2pt},
  mout/.style={rectangle,draw=FigureOutput!70,thick,fill=FigureOutput!14,align=center,
    text=BookInk,rounded corners=1.4pt,minimum height=11mm,text width=20mm,inner sep=2pt},
  mflow/.style={-{Stealth[length=1.8mm,width=1.2mm]},draw=FigureOutput!85,line width=.85pt},
  mwrite/.style={-{Stealth[length=1.8mm,width=1.2mm]},draw=FigureLoss!85,line width=.85pt},
  mlab/.style={font=\figurefont\scriptsize,text=BookMuted,fill=white,inner sep=1pt}
]
  % 上行：新行为 → 控制器
  \node[mbeh] (beh) at (12,26) {新行为 $e_k$};
  \node[mctl] (ctl) at (42,26)
    {控制器\\读键 $w^{\mathrm r}_k$／写键 $w^{\mathrm w}_k$\\擦除 $\bm e_k$／添加 $\bm a_k$};

  % 记忆矩阵（与控制器留出间隙，仅含槽位说明，无叠加文字）
  \node[mmem] (mem) at (42,0) {记忆矩阵 $\bm M$\\$m$ 个槽位};
  % CTR 与记忆矩阵同基线，箭头从右侧水平引出
  \node[mout] (ctr) at (70,0) {CTR 网络\\候选在此进入};

  \draw[mflow] (beh.east) -- (ctl.west);
  \draw[mwrite] (ctl.south) -- (mem.north);
  \node[mlab,align=left] at (49,15.5) {擦除＋写入};
  \draw[mflow] (mem.east) -- (ctr.west);

  % 读出标注独立置于记忆矩阵下方，不压框内文字、不压箭头
  \node[mlab,align=center] at (42,-18) {读出 $\bm r_k$ 与归纳状态，进入 CTR};
\end{tikzpicture}
